\subsection*{Reading the referent of a pictured expression}
When a Chinese paragraph introduced a garden and a tower and then gave the admission price of \emph{the former} and \emph{the latter}, the name-conditioned writer assigned the garden's 12 yuan to both. We first tested a free antecedent locator in human fiction. On marked pronouns in unseen works it returned a name in the correct annotated cluster on 43/200 Chinese and 157/430 English cases. Giving a second model pictures of candidate names and pooling repeated occurrences made the choice easier. With mixed-corpus fitting, the optical model answered 948/1,465 screened English GAP cases correctly, compared with 831 for the same architecture trained only on stories, 647 for the nearest name, 733 for a learned geometry control and 947 for matched byte input. On 298 further article-grouped English cases the mixed optical and byte results were 175 and 190. On Chinese CLUE the optical model accepted 54/111 genuine links and rejected 125/193 false ones. Always rejecting gave 193/304, exceeding the learned total of 179. The two scripts share a visual interface but these tasks and their class frequencies differ.

We then fitted an ordinal reference residual on paired pictured expressions. In the first sixteen occurrences screened in each of three natural reading corpora, 41 met the candidate and length screen, including adjectival uses with neither name as antecedent. The frozen base chose correctly on 14. Adding the optical residual gave 38, erasing its marked expression before training gave 13 and the matched byte residual gave 39 (Fig.~\ref{fig:story-mentions}b). The three held groups had 11 English, 15 traditional Chinese and 15 simplified Chinese cases. This is an investigator-annotated stress test of a specific construction. The model received the marked reference and its candidate spellings. A post-fit lexical rule selecting the earlier or later candidate from the expression also answered 38/41 correctly. Counts on the four previous GAP and CLUE screens were unchanged by the optical residual.

Connecting the reference reader to the existing writer recovered both admission prices in the original failure (Fig.~\ref{fig:story-mentions}a). We replaced repeated venue names with ordinal references in pairs of human dialogue utterances and reversed both venue order and questioned name. Of 539 source groups, 536 fit the length limit in both orders. Across 2,144 questions, correct attached answers rose from 604 to 1,758. Requiring all four answers per group gave 0/536 and 365/536, respectively. These constructed comparisons reuse source utterances. The difference includes clause restriction and supplied candidate names. It cannot be assigned to neural learning alone. The combined inference path contains 567,759 parameters. Inferring a cast, handling implicit referents and writing freely about new topics remain open.

The order test leaves open whether the reader can reverse its decision when a different word changes what the pronoun means. In 120 eligible contrast pairs from the Winograd Schema Collection the two names and queried pronoun stayed fixed, but the correct antecedent changed. Across the 240 variants the frozen optical reader chose correctly on 117 and a nearest-name rule on 120. The model changed its answer on four pairs and answered both variants correctly on just one. The ordering and fee results therefore coexist with a substantial failure of contextual interpretation outside the trained construction.
